\chapter{春秋：王室东迁后的诸侯竞争}
\section{时代结构}
东周王室保留礼仪与册命的象征中心，却不再拥有足以约束大国的资源。诸侯会盟、霸主调停、卿族坐大和边缘族群互动，共同构成春秋政治。

\section{编年事件}
\begin{event}{公元前770年}{平王东迁洛邑}{可靠纪年}{洛邑}
\term{经过}：平王即位并迁都洛邑，东周开始。
\term{后果}：王室依赖诸侯保护，政治重心由王室直接统治转向诸侯间竞争。
\end{event}
\begin{event}{公元前685年}{齐桓公即位，管仲相齐}{可靠纪年}{齐}
\term{经过}：齐桓公在内乱后即位，任用管仲整饬内政并组织会盟。
\term{后果}：齐成为早期霸主，说明财政、军政组织和联盟能力可以部分替代周王室的旧权威。
\end{event}
\begin{event}{公元前632年}{城濮之战}{可靠纪年}{城濮}
\term{经过}：晋、楚交战，晋文公获胜。
\term{后果}：晋取得中原盟主地位，齐、晋、楚轮番竞争的格局更加清晰。
\end{event}
\begin{event}{公元前546年}{弭兵之会}{可靠纪年}{宋}
\term{经过}：晋、楚等国在宋会盟，试图缓和长期争霸冲突。
\term{后果}：大国暂时均势并未消除，反而给各国内部卿大夫势力扩张创造空间。
\end{event}
\begin{event}{公元前481年}{春秋末年}{传统分期}{鲁}
\term{经过}：传统史学多以此年作为《春秋》记事终点。
\term{中期影响}：诸侯间的礼制秩序继续空洞化，列国开始转向更集中的军政与领土国家形态。
\end{event}
\begin{causalchain}[从霸政到变法]
霸主政治依赖联盟与礼仪承认，无法稳定解决各国内部的土地、军役和贵族分权问题。其失败为战国变法、常备军和官僚化提供了问题背景。
\end{causalchain}
